\startvolume{Paper V}{The log-Gamma phase-space atlas}

\lesson{PAPER V / V01}{Normalize the log-Gamma window}{V01_log_gamma}
{g_\alpha(y)=\Gamma(\alpha)^{-1/2}e^{\alpha y/2-e^y/2},\qquad
 \int_{\mathbb R}|g_\alpha(y)|^2dy=1}
{Read amplitude on the left and probability density on the right. A taller curve
need not have more total probability. Changing the positive shape moves the
active log position and changes its width, while the full squared integral
stays one. The curves are unrenormalized crops: the shape $0.5$ window has a
visible left tail beyond the page.}
{For $\alpha=1.5$, substitute $u=e^y$. The norm becomes
$\Gamma(1.5)^{-1}\int_0^\infty u^{0.5}e^{-u}du=1$.
The mode is $y=\log(1.5)$, whereas the mean is
$\psi(1.5)\approx0.036490$. Those two positions differ because the
log-Gamma law is asymmetric. Its variance is approximately $0.934802$.}
{Use the normalized window as the common input to every later transform.}
{V Definition 2.1 and Proposition 2.2, source lines 339 and 369. Shapes
$0.5,1.5,4$, $y\in[-7,3]$. Exact Gamma-CDF crop probabilities are recorded in
the evidence ledger; the displayed crop is not the full norm.}

\lesson{PAPER V / V02}{Build an ambiguity integral}{V02_ambiguity_integral}
{A_\alpha(h,\eta)=\int g_\alpha(y+h/2)g_\alpha(y-h/2)e^{i\eta y}dy
 =\frac{\Gamma(\alpha+i\eta)}{\Gamma(\alpha)}
 \operatorname{sech}^{\alpha+i\eta}(h/2)}
{First multiply the shifted windows. Their product localizes the integral
before modulation is applied. The second panel separates the cosine and sine
contributions to its complex value. Negative real or imaginary contributions
are phases of the integrand, not negative probability weights. Integrating both
components yields one centered ambiguity value.}
{Take $\alpha=1.5$, $h=1.4$, and $\eta=0.8$. The closed formula gives
approximately $0.536116-0.048326i$. Direct quadrature of the displayed
definition agrees. At zero modulation the sine contribution vanishes; at zero
displacement the two windows coincide and the remaining integral is the
log-Gamma characteristic function. Neither operation changes the Fourier
convention.}
{Translation and modulation are joined by one precisely centered integral.}
{V Theorem 3.1, ``Exact ambiguity formula,'' source line 404. Quadrature uses
$y\in[-30,6]$, tolerance $2\times10^{-13}$; analytic omitted-amplitude bounds
and numerical error estimates are recorded separately.}

\lesson{PAPER V / V03}{Separate magnitude from wrapped phase}{V03_phase_space}
{A_\alpha(h,\eta)=\frac{\Gamma(\alpha+i\eta)}{\Gamma(\alpha)}
 e^{-(\alpha+i\eta)\log\cosh(h/2)}\ne0
 \quad(h,\eta\in\mathbb R)}
{The left panel measures overlap strength. The right panel records the
argument modulo $2\pi$, so its sharp color transitions may be phase wraps.
They do not indicate a zero or discontinuity of the complex function.
Increasing displacement suppresses magnitude through the hyperbolic factor.
Increasing modulation introduces the separate Gamma envelope and its phase.}
{At the origin the value is exactly one. At $h=\log4$ and $\eta=0$, its
value is $(4/5)^{3/2}\approx0.715542$. At nonzero modulation, use the same
formula with a complex exponent. The real logarithm of $\cosh(h/2)$ fixes
that exponent without a branch ambiguity on the real plane. The Gamma
function has no zeros, which supplies the nonvanishing proof.}
{A phase wrap is a plotting convention; a zero is a mathematical event.}
{V Theorem 3.1 and Corollary 3.2, source lines 404 and 447. Shape $1.5$;
101-by-101 samples on $[-5,5]\times[-4,4]$. A positive sampled minimum is
finite evidence, not the universal nonvanishing argument.}

\lesson{PAPER V / V04}{Raise shape and read its differential equation}{V04_shift_differential}
{\begin{aligned}
 A_{\alpha+1}&=\frac{\alpha+i\eta}{\alpha}
 \operatorname{sech}(h/2)A_\alpha,\qquad \nu=\alpha+i\eta,\\
 [-4\partial_h^2-\nu(\nu+1)\operatorname{sech}^2(h/2)]A_\alpha
 &=-\nu^2A_\alpha.
 \end{aligned}}
{The first panel compares two constructions of the raised shape: direct
evaluation and the Gamma recurrence. The second shows the potential coefficient
in the differential equation. Both its real and imaginary parts matter when
modulation is nonzero. This complex equation should not be silently read as a
self-adjoint Hamiltonian with a real energy.}
{For $\alpha=1.5$ and $\eta=0.8$, raising shape multiplies the kernel by
$(1.5+0.8i)\operatorname{sech}(h/2)/1.5$. At $h=1.4$, the result is
approximately $0.447661+0.189299i$. Differentiate
$\operatorname{sech}^{\nu}(h/2)$ twice to obtain the displayed equation.
The Gamma prefactor is independent of $h$, so it survives unchanged in
this derivative calculation.}
{One recurrence and one differential equation describe the same kernel.}
{V Proposition 3.3 and Theorem 3.4, source lines 466 and 491.
$h\in[-5,5]$; derivatives are evaluated analytically. The residual checks
floating-point evaluation of the identity, not a discretized spectral problem.}

\lesson{PAPER V / V05}{Transform the window into a Gamma envelope}{V05_fourier_gamma}
{\widehat g_\alpha(\omega)=
 \frac{2^{\alpha/2-i\omega}\Gamma(\alpha/2-i\omega)}
 {\sqrt{2\pi\Gamma(\alpha)}},\qquad
 \rho_\alpha(\omega)=|\widehat g_\alpha(\omega)|^2}
{Keep the unitary convention
$\widehat f=(2\pi)^{-1/2}\int f(y)e^{-i\omega y}dy$ beside the plot.
The Fourier amplitude has real and imaginary parts, even though the starting
window is real. Squaring its magnitude removes the carrier phase and gives
the positive Meixner density. The second panel compares these equivalent
density constructions.}
{Set $u=e^y$. The transform becomes the Gamma integral with exponent
$\alpha/2-i\omega$ and scale $1/2$. For $\alpha=1.5$, evaluating this
integral at frequencies $-2,0,2$ agrees with the closed envelope. Conjugating
the value at $2$ gives the value at $-2$, as required for a real position
window. The positive density is consequently symmetric in frequency.}
{Preserve the Fourier sign before comparing phases or polynomial bases.}
{V Theorem 4.1, source line 540. Frequency crop $[-5,5]$. Independent
quadrature uses $y\in[-30,6]$ with explicit amplitude-tail bounds. The
ledger separates the quadrature estimate from the omitted integral.}

\lesson{PAPER V / V06}{Read depth overlap as a characteristic function}{V06_meixner_overlap}
{\rho_\alpha(\omega)=\frac{2^{\alpha-1}}{\pi\Gamma(\alpha)}
 |\Gamma(\alpha/2+i\omega)|^2,\qquad
 \int\rho_\alpha(\omega)e^{ih\omega}d\omega
 =\operatorname{sech}^{\alpha}(h/2)}
{The two panels express one probability law in frequency and characteristic
coordinates. Larger shape broadens the frequency density because its variance
is $\alpha/4$. Its characteristic function then narrows in displacement.
The right curve is also the zero-modulation ambiguity slice, so it has the
same depth-overlap interpretation as the arithmetic limit.}
{Take $\alpha=1.5$ and a four-to-one depth ratio. Then $h=\log4$ and
$\cosh(h/2)=5/4$. The exact overlap is
$(4/5)^{3/2}\approx0.715542$. This is a transform identity for the
continuous window. A finite squarefree arithmetic packet approaches it under
the compact ratio and time hypotheses of V Theorem 1.1; the identity alone
does not give a finite-depth error bound.}
{The arithmetic scale overlap becomes an exact Meixner characteristic law.}
{V Corollary 4.2, source line 582. Shapes $0.5,1.5,4$. Density crop
$[-5,5]$ is not renormalized. Full-line adaptive cosine quadrature checks
$h=0,\log4$ with declared numerical tolerance.}

\lesson{PAPER V / V07}{Resolve the negative Wigner lobes}{V07_wigner_negative}
{W_\alpha(y,\omega)=\frac1{2\pi}\int
 g_\alpha(y+h/2)g_\alpha(y-h/2)e^{-i\omega h}dh
 =\frac{2e^{\alpha y}}{\pi\Gamma(\alpha)}K_{2i\omega}(e^y)}
{The left panel shows three exact position slices. The right enlarges the
small lobes around zero. A negative value is meaningful here: the Wigner
function records phase-space interference and is not a pointwise probability
density. Its frequency symmetry does not remove those negative regions.
Changing position changes where the oscillatory Bessel lobes appear.}
{For $\alpha=1.5$, $y=0$, and $\omega=2$, the exact expression is
$2K_{4i}(1)/(\pi\Gamma(1.5))\approx-0.00155215$.
The integral remains real because its sine part is odd in the displacement
variable. Negative phase-space values therefore coexist with real-valued
Wigner slices. Integrating over an entire coordinate produces positive
marginals, which is the next lesson's separate statement.}
{Wigner negativity and positive probability marginals are compatible.}
{V Theorem 5.1, source line 615. Shape $1.5$; $y=-1,0,1$ and 241
frequency samples on $[0,4]$. Modified Bessel values use 60-digit working
precision; smooth lines connect finite samples.}

\lesson{PAPER V / V08}{Project Wigner onto its probability marginals}{V08_wigner_marginals}
{\int_{\mathbb R}W_\alpha(y,\omega)d\omega=|g_\alpha(y)|^2,
 \qquad
 \int_{\mathbb R}W_\alpha(y,\omega)dy=\rho_\alpha(\omega)}
{Integrating over frequency leaves the position density; integrating over
position leaves the frequency density. The figure draws these proved exact
marginals directly. It does not integrate a cropped Wigner image, which would
discard part of an oscillatory integral. The shaded regions show only the
probability mass inside the visible windows.}
{At shape $1.5$, the position mean is
$\psi(1.5)\approx0.036490$ and the frequency mean is zero. The position
variance is $\psi_1(1.5)\approx0.934802$, while the frequency variance is
$1.5/4=0.375$. Thus the two marginal widths differ. The negative value from
the previous lesson disappears only after complete marginal integration,
not because the Wigner function has become pointwise positive.}
{A projection can be positive even when the joint phase-space representation is not.}
{V Corollary 5.2, source line 649. Position crop $[-7,3]$, frequency crop
$[-5,5]$. Gamma-CDF position mass and adaptively integrated frequency mass
are recorded; no truncated-image marginal identity is asserted.}

\lesson{PAPER V / V09}{Compare uncertainty with orbit geometry}{V09_uncertainty_metric}
{\Delta Y\Delta P=\tfrac12\sqrt{\alpha\psi_1(\alpha)}>\tfrac12,
 \qquad g_{uu}=\alpha/4,\qquad
 g_{\eta\eta}=\psi_1(\alpha)}
{Follow the uncertainty curve toward the dashed lower bound. It approaches
that bound with increasing shape but remains strictly above it at every
finite shape. The second panel reads the same position and momentum variances
as metric coefficients of a fixed-shape Weyl orbit. Translation sensitivity
increases while modulation sensitivity decreases.}
{At $\alpha=1.5$, the momentum variance is $0.375$ and the position
variance is approximately $0.934802$. Their square-root product is
$0.592073$, greater than $1/2$. The quantum Fisher coefficients in these
two orbit directions are four times the metric, namely $1.5$ and
approximately $3.739209$. These are local parameter sensitivities, not a
claim that the packet is Gaussian at finite shape.}
{Uncertainty and geometric sensitivity use the same exact packet variances.}
{V Theorem 6.1 and Proposition 6.2, source lines 678 and 748.
Shapes $0.2\leq\alpha\leq30$. This is the translation/modulation metric
at fixed shape; mixed coefficients involving changing shape are not inferred.}

\lesson{PAPER V / V10}{Create orthogonal Laguerre packets}{V10_laguerre_gram}
{\Lambda_{n,\alpha}(y)=\sqrt{\frac{n!}{\Gamma(n+\alpha)}}
 e^{\alpha y/2-e^y/2}L_n^{(\alpha-1)}(e^y),\qquad
 \int\Lambda_{n,\alpha}\Lambda_{m,\alpha}dy=\delta_{nm}}
{Each degree adds oscillations around the same envelope. A point where two
curves cross does not determine their inner product. Orthogonality is the
weighted integral over the entire position line. The second panel displays
the degree-zero through degree-four Gram matrix, computed in the transformed
Gamma coordinate rather than by integrating the cropped curves.}
{Use $u=e^y$. The inner product becomes a generalized Laguerre integral
with weight $u^{\alpha-1}e^{-u}$. At $\alpha=1.5$, products of the shown
polynomials have degree at most eight. A 48-node generalized Gauss--Laguerre
rule integrates such polynomial products exactly in mathematical arithmetic.
Its small recorded residual measures floating-point implementation, whereas
the full orthogonality theorem covers all degrees.}
{A common envelope becomes an orthogonal basis through its natural weight.}
{V Theorem 7.1, source line 803. Shape $1.5$, degrees $0$ through $4$;
position curves cropped to $[-5,3]$. The Gram uses 48 weighted nodes,
whose exact polynomial range extends through degree 95.}

\lesson{PAPER V / V11}{Keep Fourier phases in the common ladder}{V11_fourier_ladder}
{\begin{aligned}
 \widehat\Lambda_{n,\alpha}(\omega)&=i^n\widehat g_\alpha(\omega)
 p_{n,\alpha}(\omega),\\
 J_L&=2K_0-K_+-K_-,\qquad J_M=(K_++K_-)/2.
 \end{aligned}}
{The first panel retains the degree phase $i^n$. Dropping it would replace
the actual Fourier-transformed position basis with a different frequency
basis. The second compares neighbor weights of the two real polynomial
recurrences. Their coefficients come from one abstract ladder, although
their signs and scaling differ between coordinate realizations.}
{At shape $1.5$, the first raising coefficient is
$\sqrt{1.5}$. The Laguerre neighbor has its negative, while the real
Meixner polynomial recurrence uses half that magnitude. Test
$[K_-,K_+]=2K_0$ on a finite-support basis vector below the truncation
boundary. A six-dimensional cut satisfies the interior calculation but has
a top-row defect because its last raising operation was removed.}
{Coordinate changes retain degree phases and infinite-basis boundary conditions.}
{V Theorems 8.1 and 9.1, source lines 855 and 917. Shape $1.5$; four
Fourier packet curves and a six-dimensional ladder cut. The audit records
both the small interior residual and the nonzero top-boundary defect.}

\lesson{PAPER V / V12}{Extend the arithmetic depth continuously}{V12_fractional_depth}
{C_{a,\beta}(r,0)=\sum_{j\geq0}(-1)^j\binom rj
 \zeta(a+\beta j)^{-1}
 =\sum_{j\geq0}(-1)^j\binom rj
 [\zeta(a+\beta j)^{-1}-1],\quad r>0}
{The left curve embeds the integer signed orbit in a smooth positive-depth
parameter. The right measures differences from a finite $J=80$ reference for
two equivalent infinite binomial coordinates.
The original finite partial sum retains an artificial remainder from the
unit packet. Subtracting one removes that packet and accelerates the displayed
finite reconstruction. These are signed amplitudes, not positive densities.}
{Choose arithmetic parameters $a=\beta=2$ and depth $r=2.5$. The direct
sieve sums $\mu(n)n^{-2}(1-n^{-2})^{2.5}$ through $N=200000$.
For $r>0$, the full generalized binomial coefficient sum is zero, which
justifies the centered reciprocal-zeta series. At integer depths that series
terminates. Increasing its last index improves the plotted comparison but
does not constitute a universal infinite-series error certificate.}
{Continuous arithmetic depth preserves the unit-packet boundary and signed character.}
{V Theorems 10.1 and 10.2, source lines 995 and 1020. $a>1$ is required
for the stated absolute exchanges. Direct cutoff tail is at most $1/N$;
the reference is an 80-index binomial sum at 60-digit precision. The log plot
floors differences at $10^{-18}$.}
